\section{PV-FashionMNIST 数据集构造}

\subsection{从中心分类到位置可变分类}

原始 FashionMNIST 样本是单通道 $28\times28$ 图像。PV-FashionMNIST 不改变类别标签，而是先对前景进行可选旋转，再把它粘贴到 $56\times56$ 的全零画布。以居中位置为原点，$dx>0$ 表示向右移动，$dy>0$ 表示向下移动。模型输入因此统一为 $[1,56,56]$，输出仍是十类 logits。

这个构造把类别内容与空间位置分开：对同一个原始样本，只需改变 $(dx,dy)$ 或 angle，就能得到一组内容相同而扰动强度不同的测试输入。与直接使用随机增强相比，可记录的元数据使后续按坐标和角度聚合准确率成为可能。每次取样返回
\[
  (\text{image},\text{label},\text{meta}),\qquad
  \text{meta}=\{dx,dy,\text{angle},\text{scale}\}.
\]

\subsection{数据划分与随机性控制}

FashionMNIST 官方训练集包含 60,000 张图像，项目按固定随机种子 42 划分为 54,000 张训练样本和 6,000 张验证样本；官方 10,000 张测试图像只用于最终评估。划分索引固定，避免模型之间因验证样本不同而产生额外差异。

训练期随机参数由 seed、epoch 和样本索引共同生成。这样做有两个目的：同一配置重新运行时可以复现；训练样本又不会在所有 epoch 中停留在同一位置和角度。训练循环在每个 epoch 开始时更新数据集状态，而验证集和测试集固定在 epoch 0。随机性不依赖 DataLoader 的 worker 调度，因此切换单进程或多进程加载不会改变某个样本的增强参数。

\subsection{平移协议}

项目把位置变化划分为四种强度：

\begin{itemize}
  \item \textbf{Center：}固定使用 $(dx,dy)=(0,0)$，对应原始中心分布；
  \item \textbf{Small Shift：}$dx,dy$ 分别在 $[-8,8]$ 内均匀采样，主体仍大多完整可见；
  \item \textbf{Large Shift：}$dx,dy$ 分别在 $[-18,18]$ 内均匀采样，包含明显边缘位置；
  \item \textbf{Grid Shift：}$dx,dy\in\{-18,-12,-6,0,6,12,18\}$，共形成 49 个固定位置，每个位置都评估完整测试集。
\end{itemize}

28 像素前景在 56 像素画布中无裁剪移动的上限为 14 像素。因此当 $|dx|$ 或 $|dy|>14$ 时，部分前景会越过画布边界。项目没有把越界区域缩回画面，而是按真实粘贴结果裁剪。这使 Large Shift 同时包含位置分布变化和可见信息减少。它提高了极端条件的难度，也意味着后续不能把所有性能下降都简单归因于“缺少平移不变性”。

\subsection{旋转与复合增强}

在平移实验初步建立后，数据协议进一步加入旋转。Rotation 模式把前景保持在中心，在 $[-45^\circ,45^\circ]$ 连续采样角度；Shift+Rotation 则同时在 $[-18,18]$ 内采样平移，并在相同角度范围内旋转。旋转首先作用于原始 $28\times28$ 前景，采用双线性插值和零填充，之后再粘贴到大画布。该顺序避免围绕整张 $56\times56$ 画布旋转而产生与目标无关的大范围位移。

增强模型的训练数据采用 Shift+Rotation，并以 0.25 的概率在画布上执行随机擦除。缩放参数仍作为数据集接口保留，但本轮正式实验不额外混入缩放，以便把主要变化集中在平移和旋转。除了连续随机角度，测试阶段还固定评估
\[
  \{-45^\circ,-30^\circ,-15^\circ,0^\circ,15^\circ,30^\circ,45^\circ\},
\]
从而得到可解释的角度响应，而不是只保留一次随机平均。

\subsection{具体实现}

从代码实现看，PV-FashionMNIST 在 \texttt{Dataset.\_\_getitem\_\_} 中按需生成输入。核心顺序很短：先在原始前景上执行可选旋转/缩放，再把处理后的前景交给 \texttt{\_place} 粘贴到大画布；\texttt{\_place} 只复制前景与画布重叠的区域，因此越界部分自然被裁剪。

\begin{lstlisting}[style=pvpython,caption={PV-FashionMNIST 中旋转与画布粘贴的核心片段。},label={lst:pv-dataset}]
if angle != 0.0 or scale != 1.0:
    image = TF.affine(image, angle=angle, translate=[0, 0],
                      scale=scale, shear=[0.0, 0.0], fill=0.0)

canvas = torch.zeros((1, self.canvas_size, self.canvas_size),
                     dtype=image.dtype)
canvas[:, dst_top:dst_bottom, dst_left:dst_right] = \
    image[:, src_top:src_bottom, src_left:src_right]
\end{lstlisting}

\subsection{数据构造可视化}

图\ref{fig:data-preview}从同一批类别出发，依次展示原始图像、居中画布、随机平移、纯旋转和 Shift+Rotation。它直观呈现了项目真正改变的变量：前景类别保持不变，空间位置和朝向发生变化。

\begin{figure}[H]
  \centering
  \includegraphics[width=0.98\textwidth]{data_preview.png}
  \caption{原始 FashionMNIST 与 PV-FashionMNIST 的五种输入形式。每一列保持同一类别，行间只改变画布、位置和旋转策略。}
  \label{fig:data-preview}
\end{figure}

图\ref{fig:position-variation}把同一个样本放到九个典型坐标。中心位置保留完整轮廓，四个角落则出现不同方向的截断，这也是后续网格热力图出现边缘退化的重要数据原因。

\begin{figure}[H]
  \centering
  \includegraphics[width=0.62\textwidth]{position_variation_demo.png}
  \caption{同一测试样本在九个典型 $(dx,dy)$ 位置下的外观。极端位置既改变坐标，也减少了可见区域。}
  \label{fig:position-variation}
\end{figure}

图\ref{fig:angle-variation}固定位置，只改变输入角度。$\pm45^\circ$ 下仍能辨认主体，但轮廓方向和插值后的局部像素已经与中心训练样本明显不同。位置九宫格与角度序列分别隔离了两类扰动，为后续空间和角度分析提供参照。

\begin{figure}[H]
  \centering
  \includegraphics[width=0.95\textwidth]{angle_variation_demo.png}
  \caption{同一测试样本在七个固定旋转角度下的输入变化。旋转范围由 $-45^\circ$ 延伸至 $45^\circ$。}
  \label{fig:angle-variation}
\end{figure}
